\section{MCMC fitting}
\label{sec:mcmc}

The total angular power spectrum between frequencies $\nu_i$ and $\nu_j$ is modeled as

\begin{equation}
C_\ell(\nu_i, \nu_j) = C_\ell^{\rm sync}(\nu_i, \nu_j) + C_\ell^{\rm dust}(\nu_i, \nu_j) + C_\ell^{\rm sd}(\nu_i, \nu_j),
\end{equation}

where the synchrotron, dust, and synchrotron-dust cross-correlation terms are

\begin{equation}
C_\ell^{\rm sync}(\nu_i, \nu_j) = A_s \left(\frac{\ell}{\ell_{\rm ref}}\right)^{\alpha_s} 
\left(\frac{\nu_i}{\nu_{\rm ref}}\right)^{\beta_s} \left(\frac{\nu_j}{\nu_{\rm ref}}\right)^{\beta_s},
\end{equation}

\begin{equation}
C_\ell^{\rm dust}(\nu_i, \nu_j) = A_d \left(\frac{\ell}{\ell_{\rm ref}}\right)^{\alpha_d} 
S_d(\nu_i) S_d(\nu_j),
\end{equation}

\begin{equation}
\begin{split}
C_\ell^{\rm sd}(\nu_i, \nu_j) = \rho \sqrt{A_s A_d} \Big[ S_s(\nu_i) S_d(\nu_j) \\
+ S_s(\nu_j) S_d(\nu_i) \Big] \left(\frac{\ell}{\ell_{\rm ref}}\right)^{(\alpha_s + \alpha_d)/2},
\end{split}
\end{equation}

with $S_s(\nu) = (\nu / \nu_{\rm ref})^{\beta_s}$ and the modified blackbody scaling in Rayleigh-Jeans units

\begin{equation}
S_d(\nu) = \frac{B_\nu(\nu, T_d)}{B_\nu(\nu_0, T_d)} \left(\frac{\nu}{\nu_0}\right)^{\beta_d},
\end{equation}

where $B_\nu$ is the Planck function. Here, $A_s$ and $A_d$ are the amplitudes of the synchrotron and dust components at the reference multipole $\ell_{\rm ref} = 80$ and frequencies $\nu_{\rm ref} = 11.1$~GHz and $\nu_0 = 353.0$~GHz, respectively; $\alpha_s$ and $\alpha_d$ are the multipole spectral indices; $\beta_s$ and $\beta_d$ are the frequency spectral indices; $T_d$ is the dust temperature (fixed at $T_d = 19.6$~K); and $\rho$ is the synchrotron-dust correlation coefficient.

\subsection{Data Selection and Frequency Coverage}

We perform a Markov Chain Monte Carlo (MCMC) fitting procedure using data from QUIJOTE 11~GHz; WMAP 23, 33, 41, 61, and 94~GHz; and Planck 30, 44, 70, 100, 143, 217, and 353~GHz, covering the multipole range $30 < \ell < 200$. To avoid noise bias in the polarization power spectra, we exclusively use cross-spectra between different frequency bands ($i \neq j$), discarding all auto-spectra ($i = j$). This yields 78 independent cross-spectra from the 13 frequency bands, which improves the robustness of the fitted spectral indices while mitigating systematic effects from instrumental noise.

\subsection{Bayesian Inference}

We adopt a Bayesian framework to constrain the model parameters $\boldsymbol{\theta} = \{A_s, \alpha_s, \beta_s, A_d, \alpha_d, \beta_d, \rho\}$. The posterior distribution is given by

\begin{equation}
p(\boldsymbol{\theta} \mid \mathbf{D}) \propto \mathcal{L}(\mathbf{D} \mid \boldsymbol{\theta}) \, p(\boldsymbol{\theta}),
\end{equation}

where $\mathcal{L}(\mathbf{D} \mid \boldsymbol{\theta})$ is the likelihood and $p(\boldsymbol{\theta})$ is the prior. The likelihood is constructed from the observed power spectra $\hat{C}_\ell(\nu_i, \nu_j)$ and their uncertainties:

\begin{equation}
\ln \mathcal{L} = -\frac{1}{2} \sum_{i,j,\ell} \frac{[\hat{C}_\ell(\nu_i, \nu_j) - C_\ell(\nu_i, \nu_j; \boldsymbol{\theta})]^2}{\sigma^2_\ell(\nu_i, \nu_j)}.
\end{equation}

For the power-law model fitting, we adopt uniform (flat) priors for all model parameters within physically motivated bounds:

\begin{align}
    A_s &\in [0, \infty), \\
    \alpha_s &\in [-6, 0], \\
    \beta_s &\in [-6, 0], \\
    A_d &\in [0, \infty), \\
    \alpha_d &\in [-6, 0], \\
    \beta_d &\in [0, 6], \\
    \rho &\in [-1, 1],
\end{align}

For the bin-to-bin analysis, to regularize the fit and incorporate external information from the literature, we instead impose Gaussian priors on the spectral indices:

\begin{align}
\beta_s &\sim \mathcal{N}(-3.1, 0.18^2), \\
\beta_d &\sim \mathcal{N}(1.55, 0.05^2), \\
\alpha_s &\sim \mathcal{N}(-3.0, 0.30^2), \\
\alpha_d &\sim \mathcal{N}(-2.48, 0.20^2),
\end{align}
